%-------------------------
% MASTER RESUME -- source of truth. NEVER SUBMIT THIS FILE OR ITS PDF (master/*.pdf is gitignored).
% Resume in Latex
% Author : Jake Gutierrez [Customized by Raghav Senthil Kumar]
% License : MIT
%------------------------
%
% HOW THIS FILE WORKS
%  * Superset rule: every claim on any tailored/base resume must exist here. No page limit.
%  * INVENTORY.md holds the raw answers. The comments next to each entry carry the verified facts.
%  * Trailing comment on every bullet:   % core|optional | tags: a, b
%      core     = keep by default in every tailored resume
%      optional = first to cut when space is tight
%  * Tag vocabulary: backend, frontend, fullstack, ml, llm, data, search, algorithms, hardware,
%                    testing, product, leadership, entrepreneurship
%  * Comment markers:
%      FACTS = verified, not yet worked into a bullet
%      TODO  = work for a later step
%      GAPS  = still unknown -- do not invent
%      !!    = current wording is inaccurate or overstated -- fix before any tailored resume copies it
%  * Workflow: bullets get rewritten with resume-bullet-writer, then a tech-resume-optimizer pass.
%    Rewritten so far: Clusion, Truck Pricing VLM, and (for the IBM resume) Citation Chain, EcoBin bullet 1, STEM E bullets 1+3.
%    Every other bullet below is still the ORIGINAL wording.
%  * Line-fill rule: the last line of every bullet must fill 90-100% of the line width. Fix by CUTTING words (padding is the fallback).
%    Compile to a scratch dir (pdflatex -output-directory=...) so no build files land in the repo, then check fill.

\documentclass[letterpaper,10pt]{article}

% Shared LaTeX preamble for every resume (packages, margins, section style, custom commands).
% Compile from the .tex file's own directory, e.g.  cd base && pdflatex <file>.tex
% Edit formatting HERE ONCE; do not copy it back into individual resumes.

\usepackage{latexsym}
\usepackage[empty]{fullpage}
\usepackage{titlesec}
\usepackage{marvosym}
\usepackage[usenames,dvipsnames]{color}
\usepackage{verbatim}
\usepackage{enumitem}
\usepackage[hidelinks]{hyperref}
\usepackage{fancyhdr}
\usepackage[english]{babel}
\usepackage{tabularx}
\IfFileExists{glyphtounicode.tex}{\input{glyphtounicode}}{}



\pagestyle{fancy}
\fancyhf{}
\fancyfoot{}
\renewcommand{\headrulewidth}{0pt}
\renewcommand{\footrulewidth}{0pt}

% Adjust margins
\addtolength{\oddsidemargin}{-0.5in}
\addtolength{\evensidemargin}{-0.5in}
\addtolength{\textwidth}{1in}
\addtolength{\topmargin}{-.5in}
\addtolength{\textheight}{1.0in}
\setlength{\footskip}{4.08003pt}

\urlstyle{same}

\raggedbottom
\raggedright
\setlength{\tabcolsep}{0in}

% Sections formatting
\titleformat{\section}{
  \vspace{-4pt}\scshape\raggedright\large
}{}{0em}{}[\color{black}\titlerule \vspace{-5pt}]

% Ensure that generate pdf is machine readable/ATS parsable
\ifdefined\pdfgentounicode
  \pdfgentounicode=1
\fi

%-------------------------
% Custom commands
\newcommand{\resumeItem}[1]{
  \item{
    {#1}
  }
}

\newcommand{\resumeSubheading}[4]{
  \vspace{-2pt}\item
    \begin{tabular*}{0.97\textwidth}[t]{l@{\extracolsep{\fill}}r}
      \textbf{#1} & #2 \\
      \textit{#3} & \textit{#4} \\
    \end{tabular*}\vspace{-7pt}
}

\newcommand{\resumeSubSubheading}[2]{
    \item
    \begin{tabular*}{0.97\textwidth}{l@{\extracolsep{\fill}}r}
      \textit{\small#1} & \textit{\small #2} \\
    \end{tabular*}\vspace{-7pt}
}

\newcommand{\resumeProjectHeading}[2]{
    \item
    \begin{tabular*}{0.97\textwidth}{l@{\extracolsep{\fill}}r}
      #1 & #2 \\
    \end{tabular*}\vspace{-7pt}
}

\newcommand{\resumeSubItem}[1]{\resumeItem{#1}\vspace{-4pt}}

\renewcommand\labelitemii{$\vcenter{\hbox{\tiny$\bullet$}}$}

\newcommand{\resumeSubHeadingListStart}{\begin{itemize}[leftmargin=0.15in, label={}]}
\newcommand{\resumeSubHeadingListEnd}{\end{itemize}}
\newcommand{\resumeItemListStart}{\begin{itemize}[leftmargin=0.22in, itemsep=1pt, topsep=2pt, parsep=0pt, partopsep=0pt]}
\newcommand{\resumeItemListEnd}{\end{itemize}\vspace{-2pt}}


\begin{document}

\begin{center}
    {\Huge\bfseries Raghav Senthil Kumar} \\ \vspace{1pt}
    \small Phoenix, AZ $|$ 623-332-2168 $|$ \href{mailto:raghavs9@illinois.edu}{\underline{raghavs9@illinois.edu}} $|$
    \href{https://linkedin.com/in/raghav-senthil-kumar}{\underline{linkedin.com/in/raghav-senthil-kumar}} $|$
    \href{https://github.com/Raghavsk24}{\underline{github.com/raghav}}
\end{center}


%-----------EDUCATION-----------
% Graduating in 3 years (May 2029, confirmed). No college GPA yet. No high school entry.
% Coursework is exactly as listed. Activities stay one line until roles/contributions exist.
\section{Education}
  \resumeSubHeadingListStart
    \resumeSubheading
      {University of Illinois Urbana-Champaign}{Champaign, IL}
      {B.S. in Computer Science and Statistics}{Expected: May 2029}
      \resumeItemListStart
        \resumeItem{\textbf{Coursework:} Computer Science I (Java), Computer Science II (C++), Data Science \& ML (Python), Calculus III}
        \resumeItem{\textbf{Activities:} Disruption Lab, SIGRobotics, AI Alignment, BuildIllinois, Innovation LLC, Agentic AI}
      \resumeItemListEnd
  \resumeSubHeadingListEnd


%-----------EXPERIENCE-----------
\section{Experience}
  \resumeSubHeadingListStart

    % ENTRY: Clusion | tags: backend, search, fullstack | priority: high
    % FACTS: team of 10 engineers, Agile. The org built on AWS and ran CI/CD; hosting was handled by the org.
    % FACTS: each engineer owned separate features. Raghav built the ENTIRE document search engine for debate evidence.
    % FACTS: stack for that part = Python FastAPI backend, Elasticsearch. The frontend was Next.js (how much Raghav built is unclear, so it is not claimed).
    % FACTS: the IQR filter ran on the word-count distribution of tag, citation and evidence text.
    % FACTS: relevance favors recency (results skew to the current month's topic, e.g. October 2025 results in October 2025).
    % FACTS: "10K+" = people who have visited the site, NOT registered or active users. Always say "visitors".
    % FACTS: revenue "couple thousand" (currency and period unconfirmed). NOT used in any bullet.
    % DONE (bullet-writer + tech optimizer): ownership fixed ("Built a SaaS platform" -> built the search engine), users -> visitors,
    %       AWS / Agile / team of 10 back-ported with accurate ownership (the GoDaddy tailored bullet still has the old overstated wording).
    % GAPS: how recency was implemented (Elasticsearch date boost/decay?); retrieval-quality metric; which AWS services the org used; revenue figure.
    \resumeProjectHeading
      {\textbf{Founding Software Engineer $|$ Clusion}}{September 2025 -- January 2026}
      \resumeItemListStart
        \resumeItem{Built the document search engine (\emph{Python FastAPI}, \emph{Elasticsearch}) for a competitive high school debate research SaaS platform hosted on \emph{AWS}, developed by an \emph{Agile} team of \textbf{10 engineers} and drawing \textbf{10K+ site visitors}} % core | tags: backend, search, product
        \resumeItem{Indexed \textbf{600K+ debate-evidence documents} on \emph{Elasticsearch} via \emph{FastAPI}, returning results in \textbf{under 1 second}} % core | tags: backend, search
        \resumeItem{Designed and built a text-retrieval algorithm that extracts evidence segments from debate cases by validating source URLs and filtering statistical outliers with \textbf{IQR} on the word-count distribution of tag, citation and evidence text} % core | tags: backend, algorithms, data
        \resumeItem{Developed a de-duplication algorithm that flagged \textbf{30\%} of evidence documents as identical, reducing index bloat} % optional | tags: backend, algorithms, data
        \resumeItem{Implemented recency-based relevance so evidence from the current month's debate topic ranks first in results} % optional | tags: backend, search
      \resumeItemListEnd

    % ENTRY: Jetson | tags: product, data | priority: low for SWE, high for product roles
    % FACTS: ed-tech company that teaches entrepreneurship.
    % FACTS: 3 features tested = user onboarding, gamified leaderboard, community hub.
    % FACTS: onboarding A/B test ran one week and all 30 participants completed it. The onboarding redesign shipped.
    % GAPS: what the report recommended beyond onboarding; completion rate of the old flow (needed for a true before/after).
    \resumeProjectHeading
      {\textbf{Product Development Intern $|$ Jetson}}{January 2025 -- February 2025}
      \resumeItemListStart
        \resumeItem{Conducted \textbf{7 iterations of UI/UX testing} across 3 experimental features on the Jetson App (100K+ users)} % optional | tags: product
        \resumeItem{Delivered a \textbf{5-page report} recommending strategies to optimize user engagement directly to Jetson's CEO} % optional | tags: product, leadership
        \resumeItem{Redesigned user onboarding on \emph{Figma}, cutting it from \textbf{12 to 5 pages}, validated via \textbf{A/B testing} on 30 test users} % core | tags: product, frontend
      \resumeItemListEnd

    % ENTRY: STEM E | tags: data | priority: low for SWE
    % FACTS: STEM-awareness nonprofit. Tools: Hootsuite, Power BI, Microsoft Excel, plus some Python and pandas.
    % FACTS: pitched to the Executive Director, who adopted the strategy. The 24% lift was measured over the following three-month quarter.
    % DONE: bullet 1 now lists Python (pandas) and Power BI alongside Hootsuite and Excel (all confirmed used).
    % DONE: bullet 3 states the Executive Director adopted the strategy. CUT FOR LINE-FILL (still true): the 24% was measured over the following quarter.
    % GAPS: the worksheet says "how well the reviews were" -- confirm the metric is views/engagement before it goes in a bullet.
    \resumeProjectHeading
      {\textbf{Data Analytics Intern $|$ STEM E}}{June 2024 -- July 2024}
      \resumeItemListStart
        \resumeItem{Analyzed 3 years of TikTok data (\textbf{200+ posts}) with \emph{Python} (\emph{pandas}), \emph{Hootsuite}, \emph{Power BI} and \emph{Microsoft Excel}} % optional | tags: data
        \resumeItem{Identified that ``science fun facts'' posts comprised 69.3\% of views, despite accounting for 7\% of all content} % optional | tags: data
        \resumeItem{Pitched a ``fun facts'' strategy to the Executive Director, who adopted it; engagement rose \textbf{24\% (8K+ views)}} % core | tags: data, product
      \resumeItemListEnd

  \resumeSubHeadingListEnd


%-----------PROJECTS-----------
% Rendered Stack lines are deferred to the tech-resume-optimizer pass (keyword decision).
\section{Projects}
    \resumeSubHeadingListStart

    % ENTRY: Truck Pricing VLM | tags: ml, llm, data, leadership | priority: high
    % FACTS: led a team of 4 and designed the entire algorithm. Built for Kamion (YC S22); no Kamion feedback. Result: 2nd of 7 teams at the UIUC Founder's Hackathon.
    % FACTS: dataset = 15,000 scraped listing images. Audit = de-duplicated images and removed non-USD price listings.
    % FACTS: after every wrong prediction against labeled listings the JSON schema was revised, until the model labeled 10 consecutive images correctly for each of 3 features (condition, brand, primary subject).
    % FACTS: the final schema extracts condition, brand, model and year, and filters non-truck images.
    % FACTS: 87% accuracy = 87% of 30 samples (about 26/30).
    % DONE (bullet-writer + tech optimizer): "RLHF"/"fine-tuned" removed (it was schema revisions). Sample size (30) now stated next to the 87%.
    % !! Never say RLHF or fine-tuned.
    % GAPS: (1) the "30": 30 schema revisions, or 10 images x 3 features = 30 correct labels? Bullets state only the 10-in-a-row criterion until confirmed.
    %       (2) what "87% accuracy" means for a price (within +/-X% of listing price?) and whether the 30 samples were held out.
    \resumeProjectHeading
          {\textbf{Truck Pricing VLM} $|$ \href{https://truck-pricing-vlm.vercel.app/}{\underline{Live Website}} $|$ \href{https://github.com/Raghavsk24/Truck-Pricing-VLM}{\underline{Source Code}}}{September 2026}
          \resumeItemListStart
            \resumeItem{Led team of \textbf{4} to build a pricing algorithm for \textbf{Kamion (YC S22)}. \textbf{\underline{2nd of 7 at UIUC Founder's Hackathon}}} % core | tags: ml, llm, leadership
            \resumeItem{Scraped \textbf{15,000} truck-listing images for \emph{data mining}, auditing quality via de-duplication and non-USD filtering} % optional | tags: data, ml
            \resumeItem{Designed a \emph{JSON schema} for a \textbf{Claude Sonnet 4.6} \emph{agent} to extract condition, brand, model and year from photos} % core | tags: ml, llm
            \resumeItem{Revised the schema after each wrong prediction until \textbf{10 consecutive images} were correct on each of \textbf{3 features}} % core | tags: ml, llm, testing
            \resumeItem{Built a \emph{quantile regression} pricing model on the extracted features, reaching \textbf{87\% accuracy} on \textbf{30 samples}} % core | tags: ml, algorithms
          \resumeItemListEnd

    % ENTRY: SchoolPulse AI | tags: ml, backend, hardware, leadership | priority: medium
    % FACTS: led a team of 5. Finalist at the USAII Global Hackathon (10 teams -- confirm whether 10 finalists or 10 total). Voice LLM = Gemma 4.
    % FACTS: tools named so far = AI trash can, water level sensor.
    % GAPS: name of the 3rd tool; data ingestion rate; anomaly-detection results; whether it ran at a real school.
    \resumeProjectHeading
          {\textbf{SchoolPulse AI} $|$ \href{https://school-pulse-ai.vercel.app/}{\underline{Live Website}} $|$ \href{https://github.com/vignesh-nagarajan-vn/SchoolPulse-AI}{\underline{Source Code}}}{June 2026}
          \resumeItemListStart
            \resumeItem{Engineered a \textbf{suite of 3 edge AI tools} that collects data regarding a school's energy, water and food consumption} % optional | tags: ml, hardware
            \resumeItem{Built a \emph{Python FastAPI} backend that ingests live water, waste \& energy data from all 3 tools, runs \textbf{4 Random Forest models} to detect anomalies \& deliver actionable steps via \textbf{\emph{Gemma} Voice LLM} on a \emph{Next.js} dashboard} % core | tags: backend, ml, llm, fullstack
            \resumeItem{Led a team of 5 across project ideation, development and presentation. \textbf{\underline{Finalist at USAII Global Hackathon}}} % core | tags: leadership
          \resumeItemListEnd

    % ENTRY: Citation Chain | tags: frontend, backend, llm, data | priority: high
    % FACTS: paper data is stored on Supabase (PostgreSQL). Raghav wrote SQL to filter papers by field, year and citation count.
    % DONE: the old run-on bullet is split into ingestion and RAG bullets; a SQL bullet was added.
    % CUT FOR LINE-FILL (still true, re-add if room): ingestion is in batched requests of 100; RAG answers questions about the citation graph via the Vercel AI SDK.
    % GAPS (all still ??? in INVENTORY.md): graph size and depth, load times, what gets embedded and how it is chunked,
    %       caching / rate limits, TypeScript vs JavaScript, solo vs team, users.
    \resumeProjectHeading
          {\textbf{Citation Chain} $|$ \href{https://citation-chain-two.vercel.app/}{\underline{Live Website}} $|$ \href{https://github.com/Raghavsk24/Citation-Chain}{\underline{Source Code}}}{June 2026}
          \resumeItemListStart
            \resumeItem{Built an interactive citation graph (\emph{React Flow}) that maps a paper's outbound citations across \textbf{26 research fields}} % core | tags: frontend, data
            \resumeItem{Built a dual-API ingestion pipeline resolving papers via \emph{OpenAlex} \& backfilling metadata from \emph{Semantic Scholar}} % core | tags: backend, data
            \resumeItem{Deployed a \textbf{RAG} \emph{agent} (\emph{Claude Sonnet 4.5}) that retrieves session-scoped vector documents (\emph{Supermemory})} % core | tags: backend, llm, data
            \resumeItem{Stored paper data on \emph{Supabase} (\emph{PostgreSQL}) and wrote \emph{SQL} queries to filter papers by field, year and citation count} % core | tags: backend, data
          \resumeItemListEnd

    % ENTRY: EcoBin | tags: ml, hardware, data | priority: high
    % FACTS: paper on arXiv (2606.15547), 2 datasets on Kaggle, 2nd at the Arizona State Science Fair.
    % FACTS: framework = PyTorch (confirmed by Raghav). NumPy was also used in EcoBin (where exactly is not stated, so NumPy stays a skills-line item).
    % GAPS (all still ??? in INVENTORY.md): contamination-detection F1, test-set size, whether TensorFlow/Keras were also used,
    %       what ran inference and how the 2 s latency was measured, Kaggle dataset names, solo vs team.
    \resumeProjectHeading
          {\textbf{EcoBin} $|$ \href{https://arxiv.org/abs/2606.15547}{\underline{Research Paper}} $|$ \href{https://github.com/Raghavsk24/EcoBin}{\underline{Source Code}}}{March 2026 -- May 2026}
          \resumeItemListStart
            \resumeItem{Designed a two-stage \emph{EfficientNetV2-S} neural network (\textbf{4.39M} parameters) in \emph{PyTorch}, trained on \textbf{15K+ images} to classify waste into 4 disposal pathways at \textbf{96.13\% accuracy}. Published a paper on \emph{arXiv} and 2 datasets on \emph{Kaggle}} % core | tags: ml, data
            \resumeItem{Built a synthetic data generation pipeline using \emph{U2-Net} segmentation and alpha-compositing to generate 9K images to train the model to detect contamination in recyclable waste. \textbf{\underline{Placed 2nd at the Arizona State Science Fair}}} % core | tags: ml, data
            \resumeItem{Engineered a smart bin prototype using \emph{Arduino Uno} that sorts the waste item with \textbf{less than 2 second latency}} % optional | tags: hardware
          \resumeItemListEnd

    % ENTRY: Chronos | tags: fullstack, algorithms, entrepreneurship | priority: high
    % FACTS: won the 2025 Congressional App Challenge. Licensed to 8 startups, 2 LOIs worth $3.5K.
    % GAPS (all still ??? in INVENTORY.md): how slots are ranked, usage numbers, whether 30+ people was tested or is a design limit,
    %       what "licensed" meant and the LOI terms, APIs used, Congressional App Challenge level, solo vs team.
    \resumeProjectHeading
          {\textbf{Chronos} $|$ \href{https://usechronos.live/}{\underline{Live Website}} $|$ \href{https://github.com/Raghavsk24/Chronos}{\underline{Source Code}}}{December 2025}
          \resumeItemListStart
            \resumeItem{Built a \textbf{scheduling agent} (\emph{React/Firebase}) that auto-books meetings for large groups (30+ ppl) on Google Calendar} % core | tags: fullstack
            \resumeItem{Designed a \textbf{5-step \emph{Python} scheduling algorithm} that scans calendar availability, returns a ranked list of time slots and sends email confirmation for each booking. Validated algorithm across \textbf{14 unit and integration test cases}} % core | tags: backend, algorithms, testing
            \resumeItem{Created a lobby system that enables the host to authenticate invited participants and queue follow-up meetings} % optional | tags: fullstack
            \resumeItem{Licensed Chronos to \textbf{8 startups} \& signed \textbf{2 LOIs worth \$3.5K}. \textbf{\underline{Won 2025 Congressional App Challenge}}} % core | tags: entrepreneurship, product
          \resumeItemListEnd

    \resumeSubHeadingListEnd


%-----------TECHNICAL SKILLS-----------
% Superset of every skill on any resume (base, GoDaddy, State Farm) plus pandas/Power BI/Hootsuite/Excel from STEM E.
% Evidence map = INVENTORY.md section 4. Confirmed true but no evidence written yet: Node.js, MySQL, Git CI/CD, API Development, NumPy (used in EcoBin, placement unstated).
% Evidenced 2026-09-20: SQL, PostgreSQL, Supabase (Citation Chain bullets); PyTorch (EcoBin bullet 1).
% No evidence written yet (defend or remove): Flask, HTML/CSS, JavaScript, Hugging Face, Google Cloud, Docker, Bun,
%   Claude Code, Codex, TensorFlow, Keras, scikit-learn, OpenCV, Grad-CAM.
% TypeScript is NOT listed: usage unknown.
% KEYWORD TERMS added 2026-09-20 for the five IBM resumes (Raghav: "add as many relevant keywords as possible"). Each is a description of work
% already in the bullets above. Be ready to explain every one in an interview:
%   Data Pipelines / Data Ingestion / ETL = Citation Chain (pulls papers from OpenAlex + Semantic Scholar, backfills metadata, stores on Supabase).
%   Data Quality / Validation / Cleaning  = Clusion (de-duplication, source-URL validation, IQR outlier filter); Truck VLM (de-dup, non-USD filter).
%   Structured / unstructured data        = Truck VLM (photos -> structured features via JSON schema); Clusion (debate-evidence text).
%   Data Mining (Raghav's term)           = Truck VLM scraping of 15,000 listing images.
%   AI Agents (Raghav's term)             = Truck VLM Claude agent, Citation Chain RAG agent, Chronos scheduling agent.
%   Prompt Engineering                    = Truck VLM JSON schema revised after each wrong prediction.
%   REST APIs                             = FastAPI backends; consumed OpenAlex, Semantic Scholar, Google Calendar APIs (confirm "REST" if asked).
%   Full-Stack Development                = Chronos (React/Firebase), SchoolPulse (FastAPI + Next.js), Clusion.
%   Unit Testing / Authentication         = Chronos (14 unit and integration tests; lobby authenticates invited participants).
%   Vector Store                          = Supermemory holds Citation Chain's session-scoped vector documents.
%   Microsoft: Power BI and Excel only (STEM E). Not Azure, .NET, C#, Power Platform apps or ServiceNow.
% NOT claimed anywhere (no evidence): Azure, .NET, C#, ServiceNow, Kubernetes, Spark, Kafka, Airflow, dbt, Databricks, Snowflake, Scala, R,
%   TypeScript, Go, LangChain, LangGraph, LlamaIndex, MCP, Cursor, GitHub Copilot, Tableau, embeddings work, data modeling, governance.
% Category names/grouping get finalized in the tech-resume-optimizer pass.
\section{Technical Skills}
 \begin{itemize}[leftmargin=0.15in, label={}]
    \item{
     \textbf{Languages}{: Java, Python, C++, JavaScript, SQL, HTML/CSS} \\
     \textbf{Frameworks \& Libraries}{: FastAPI, Flask, Next.js, React, React Flow, Node.js} \\
     \textbf{Databases \& Search}{: Elasticsearch, PostgreSQL, MySQL, Firebase, Supabase} \\
     \textbf{Cloud \& DevOps}{: AWS, Google Cloud, Vercel, Docker, Git, GitHub, CI/CD} \\
     \textbf{AI/ML}{: PyTorch, TensorFlow, Keras, scikit-learn, OpenCV, Grad-CAM, Hugging Face, Claude API, Gemma 4, Vercel AI SDK, Supermemory, RAG} \\
     \textbf{Data \& Analytics}{: pandas, NumPy, Power BI, Microsoft Excel, Hootsuite, A/B Testing} \\
     \textbf{Tools \& Practices}{: Agile, API Development, Google Calendar API, OpenAlex API, Semantic Scholar API, Figma, Kaggle, Arduino IDE, Bun, VS Code, Claude Code, Codex} \\
     \textbf{Concepts}{: Data Pipelines, Data Ingestion, ETL, Data Quality, Data Validation, Data Cleaning, Data Mining, Data Visualization, Statistical Analysis, Prompt Engineering, AI Agents, LLMs, Generative AI, REST APIs, Full-Stack Development, Unit Testing, Authentication}
    }
 \end{itemize}


%-----------CERTIFICATIONS-----------
% Both are from INVENTORY.md. Keep this section only for resumes with room.
\section{Certifications}
 \begin{itemize}[leftmargin=0.15in, label={}]
    \item{
     \begin{tabular*}{0.97\textwidth}[t]{l@{\extracolsep{\fill}}r}
       % core
       \textbf{Harvard CS50x:} Introduction to Computer Science & November 2024 \\
       % optional -- "Basic" level reads weak, so cut this one first
       \textbf{HackerRank Problem Solving (Basic)} Certificate & August 2026 \\
     \end{tabular*}
    }
 \end{itemize}
 \vspace{3.53pt}

%-------------------------------------------
\end{document}
